\scanpage{29}
Comparons cette équation fonctionnelle avec celle pour la fonction $L_E$, on trouve, en \deleted{\unsure{substituant}} posant pour simplifier $p^{n+\rho}=q$, et $A(t)\delta(E)(-t)^{\chi(E)}=B(t)$ :
\[
 L\left(\frac{1}{qt}\right)=B(t)L(t)
\]
et faisant $t\rightsquigarrow\frac{1}{qt}$ \unsure{dedans}, on trouve
\[
 L(t)=B\left(\frac{1}{qt}\right)L\left(\frac{1}{qt}\right)
      =B\left(\frac{1}{qt}\right)B(t)L(t)
\]
d'où
\[
 B(t)B\left(\frac{1}{qt}\right)=1\quad .\quad .
\]
Remplaçons-y $B$ par sa valeur, on trouve
\[
 L_E^\infty\left(\frac{1}{\deleted{p}qt}\right)^{-1}
 \delta(E)(-t)^{\chi(E)}L_E^\infty(t)^{-1}\delta(E)
 \left(-\frac{1}{qt}\right)^{\chi(E)}=1
\]
i.e.
\[
 \deleted{L_E^\infty\left(\frac{1}{p^{n+\rho}t}\right)^{\unsure{1|-1}}}
 \quad L_E^\infty(t)L_E^\infty\left(\frac{1}{qt}\right)
 =\delta'(E)^2\deleted{\unsure{\frac{1}{qt}|\frac{t}{qt}}}q^{-\chi(E)}
\]

C'est précisément l'équation fonctionnelle \unsure{explicitée|explicité} plus haut.

Posons
\[
 \unsure{\xi|\zeta}_E(t)=t^{\frac{\chi(E)}{2}}L_E(t)
\]
alors
\[
 \begin{aligned}
 \xi_E\left(\frac{1}{qt}\right)&=\left(\frac{1}{qt}\right)^{\chi(E)/2}
 L_E\left(\frac{1}{qt}\right)=\left(\frac{1}{qt}\right)^{\chi(E)/2}
 A(t)\delta(E)(-t)^{\chi(E)}L_E(t)\\
 &=\underbrace{q^{-\chi(E)/2}A(t)\delta(E)}\,(-1)^{\chi(E)}
 \underbrace{t^{\chi(E)/2}L_E(t)}_{\xi_E(t)}
 \end{aligned}
\]
i.e.
% archive sheet number: 28

